%*****************************************
\cleardoublepage
\phantomsection
\pdfbookmark[1]{Appendix}{appendix}
\chapter*{Appendix}\label{ch:appendix}
\markboth{\spacedlowsmallcaps{Appendix}}{\spacedlowsmallcaps{Appendix}}
\addcontentsline{toc}{chapter}{Appendix}
%*****************************************

This appendix provides detailed technical specifications, configurations, procedural descriptions, and relevant code segments.
Those details support the reproducibility by providing technical insights at relevant points.

\section*{Prompt Templates and Code}\label{app:template_code}
The complete prompt templates and the code used are available in a GitHub repository at:\\

\url{https://github.com/HSKuenzig/bachelor-thesis}

The question sets, and the generated data sets are provided at a dedicated repository available with the necessary access permissions.
%
\section*{Chunking Strategy Details and Implementation}\label{app:chunking_strategies}

This section extents the description in\cref{sub:hybrid_chunking_sol}.
It provides a more detailed configuration setup and concrete code listings of the two chunking strategies implemented within the application's hybrid chunking strategy.

\subsection*{1. Semantic Chunking}
For the initial segmentation of the original text of the knowledge base, a semantic approach is adopted - \textbf{interquartile chunking}. 
This approach calculates the semantic distance between consecutive sentences.
A new semantic section is identified as soon as the shift in context exceeds the calculated interquartile distance threshold.
Major changes in topic are thus effectively treated as statistical outliers.

Therefore, \textit{LangChain}'s \texttt{interquartile} method is chosen and implemented by modifying the original template from \cite{langchainsemantic}.\\
The created \textit{LangChain} \texttt{SemanticChunker} instance is configured with \texttt{FastEmbedEmbeddings} using the default \texttt{BAAI/bge-small-en} embedding model (see \cref{par:api_generation_rationale}).
The \texttt{breakpoint\_threshold\_type} is set to \texttt{"interquartile"}.

The interquartile range can be scaled by a keyword argument:
\begin{itemize}
    \item \texttt{breakpoint\_threshold\_amount}: A float value that scales the interquartile range.
    By adjusting this value, the sensitivity of the chunking process can be controlled.
\end{itemize}
The respective threshold was unmodified and uses the default of $1.5$.
%
Within the \ac{rag} pipeline the source text is loaded, \textit{semantic chunks} are generated.
The generated chunks are further processed and their metadata is expanded with relevant entries (e.g., \texttt{chunk\_type="semantic"}, \texttt{chunk\_id}, \texttt{total\_chunks}) for subsequent evaluation.
The resulting metadata allows for later filtering and further precise processing of relevant chunks.

\subsection*{2. Fixed-Size Chunking}\label{app:sub:fixed_size_chunking}
The secondary and added path provides the fallback to retrieve relevant input as context for the generation step.
%
\textbf{Fixed-Size Chunking} does not focus on the preservation of contextual relationships, as described in \cref{sub:chunking_prel}.
Nevertheless, this method provides an artificial contextual relationship between the chunks due to content overlap.

A \texttt{RecursiveCharacterTextSplitter} is instantiated from \textit{LangChain}.
The \textit{splitter} uses the \texttt{cl100k\_base} tokeniser provided by the \texttt{tiktoken} library.
As a result, the knowledge base is first tokenised.
On this basis, the chunk length can be determined using the \textit{token} count rather than the raw characters.
The realistic token count is relevant for subsequent \ac{api} interaction in the process, as it allows for the concrete control and calculation of the \textit{context window}.

Following the \textit{tokenisation} the corresponding \textit{splitter} creates chunks with:
\begin{itemize}
    \item \texttt{chunk\_size\_tokens}: The maximum number of tokens allowed in a single chunk - Set to $1000$
    \item \texttt{chunk\_overlap\_tokens}: The number of tokens that overlap between consecutive chunks - Set to $400$
\end{itemize}
These chunks were tagged with multiple metadata entries relevant to subsequent evaluation (see \cref{lst:output_file_structure_answered_rag_para})

This particular setting of the chunk size and contextual overlap is selected to ensure preservation of semantic context across chunk boundaries without creating excessive redundancy.

\section*{Vector Storage and Retrieval}\label{app:vector_storage_sretrieval}
This section supplements the description of the storage and indexing workflow outlined in \cref{subsub:storage_retrieval_db}.

Data ingestion, persistence, and vector-similarity queries are managed by a centralized abstraction class, the \texttt{QdrantManager}.
This orchestrating \enquote{manager} maps application operations to network calls directed at the underlying vector database.
Each \textit{chunk} (whether semantic or fixed-size) generated by the process of \textit{chunking} is converted into a structured \ac{db} record.
Such a \ac{db} record is referred to as a \textit{point}.
A standard \textit{Qdrant point} is defined by three primary elements:
\begin{enumerate}
    \item \textbf{ID}: A globally unique identifier (UUID v4 or integer) assigned to every chunk.
    These \textit{ID}s allow for precise filtering and analysis of the stored data.
    \item \textbf{Vector}: The mathematical representation (embedding) of the text chunk (see \cref{subsub:vocab_mapping}).
    As implemented, the embedding process produces vectors with $384$ dimensions using cosine similarity (see \cref{sec:vector_simil_search}).
    \item \textbf{Payload}: The actual text content.
    As implemented, the payloads contain the text chunks from \citet{bb2} combined with all parsed metadata.
    The metadata holds values such as the \texttt{source} and \texttt{chunk\_type} (e.g, \texttt{"semantic"} or \texttt{"fixed-size"}).
\end{enumerate}

The database engine runs isolated in a \textit{Docker container} using the official \textit{Qdrant} registry image which is automatically pulled on initial container startup.
The resulting collection holding all \texttt{points} is mounted to a persistent \textit{Docker volume}.
This setup ensures system persistence and prevents overhead costs on restart.
Overall, the configuration is designed to provide reproducibility and consistent storage of the vectorised and embedded knowledge base. 
%
\section*{Installation and Setup}\label{app:installation_setup}
This section describes the setup procedures and required configurations in order to deploy the containerised application environment.

\subsection*{Prerequisites}
To run the pipeline on a machine requires the following dependencies installed on the host system:
\begin{enumerate}
    \item \textbf{Git}: Version control client. Necessary to clone the code repository.
    \item \textbf{Docker Service}: Docker Desktop or Docker Engine (latest version is recommended) to orchestrate the containerised services.
    \item \textbf{API Services}: At least one \ac{api} authorisation key loaded as an environment variable:
    \begin{itemize}
        \item \texttt{OPENROUTER\_API\_KEY}: Set to access models from \textit{OpenRouter}.
        \item \texttt{BLABLADOR\_API\_KEY}: Set to access models from \textit{Blablador}.
    \end{itemize}
    If other providers are chosen, the corresponding \ac{api} key and possibly the internal model mapping must be adapted.
\end{enumerate}

\subsection*{Deployment}\label{app:deployment}
These following steps are required to deploy and launch the application:

\begin{enumerate}
    \item \textbf{Clone the Repository}:
    Clone the codebase from the primary version control host and enter the root directory:
\begin{lstlisting}[language=bash, caption={Cloning source directory}, xleftmargin=\leftmargin]
git clone https://github.com/HSKuenzig/bachelor-thesis-code
cd bachelor-thesis-code
    \end{lstlisting}

    \item \textbf{Configure Environment Variables}:
    The project workspace contains a pre-configured template file named \texttt{.env}.
    Edit this file to supply your \ac{api} authorisation keys:
\begin{lstlisting}[language=bash, caption={Environment file keys}, xleftmargin=\leftmargin]
OPENROUTER_API_KEY=your_openrouter_key
BLABLADOR_API_KEY=your_blablador_key
SITE_URL=http://localhost:8501
\end{lstlisting}

    \item \textbf{Orchestrate Containerised Services}:
    Build and launch the stack (including the \ac{db} and web client services):
\begin{lstlisting}[language=bash, caption={Starting docker service stack}, xleftmargin=\leftmargin]
docker compose up --build -d
\end{lstlisting}
    \item \textbf{Access the Application}:
    Open a web browser and navigate to the application endpoint hosted locally:
    \begin{flushleft}
    \url{http://localhost:8501}
    \end{flushleft}
\end{enumerate}
